% !TeX root = ../main.tex
% 02-introduction.tex
\chapter{Einleitung}
\label{chap:introduction}


\section{Motivation und gesellschaftlicher Kontext}
\label{sec:motivation}

\lipsum[1-3]


\section{Problemstellung}
\label{sec:problem}

\lipsum[3-9]


\section{Zielsetzung und Forschungsfragen}
\label{sec:goal_research}
\lipsum[10]

Hieraus leitet sich die folgende übergeordnete Forschungsfrage ab:

\begin{quote}
    \textit{\lipsum[1][11]}
\end{quote}

Zur strukturierten Beantwortung werden vier \acp{TF} definiert, deren thematische Einordnung sowie der jeweilige Bezug zu den \aclp{AP}n~(\acs{AP}) und den bearbeitenden Kapiteln in Tabelle~\ref{tab:forschungsfragen} zusammengefasst sind.

% !TeX root = ../main.tex

\begin{table}[htbp]
    \centering
    \scriptsize
    \caption{Teilforschungsfragen, Arbeitspakete und Einordnung in die Kapitelstruktur}
    \label{tab:forschungsfragen}
    \begin{threeparttable}
        \renewcommand{\arraystretch}{1.6}
        \begin{tabularx}{\textwidth}{c X c c}
            \toprule
            \textbf{Nr.} & \textbf{Forschungsfrage} & \textbf{AP} & \textbf{Kap.} \\
            \midrule

            \textbf{\acs{TF}\,1} & \lipsum[1][12] & 2,\;3 & 3,\;4 \\

            \textbf{\acs{TF}\,2} & \lipsum[1][13] & 1,\;4 & 3,\;5 \\

            \textbf{\acs{TF}\,3} & \lipsum[1][14] & 5,\;6 & 3,\;6 \\

            \textbf{\acs{TF}\,4} & \lipsum[1][15] & 5,\;6 & 7 \\

            \bottomrule
        \end{tabularx}
    \end{threeparttable}
\end{table}

Die abschließende Synthese aller Teilantworten sowie die Beantwortung der übergeordneten Forschungsfrage erfolgt in Kapitel~\ref{chap:conclusion_outlook}.


% End of File